\documentclass[11pt, a4paper]{article}

\usepackage[T1]{fontenc}
\usepackage[utf8]{inputenc}
\usepackage{tgpagella}
\usepackage{tgcursor} % Soluciona el warning de fuente monoespaciada en negrita
\usepackage{microtype}
\usepackage[spanish, es-noshorthands]{babel}
\usepackage{amsmath, amssymb} % Soluciona el error "Undefined control sequence \square"
\usepackage[table, xcdraw]{xcolor}
\usepackage{booktabs}
\usepackage{tabularx}
\usepackage[margin=2cm]{geometry}
\usepackage{fancyhdr}

\pagestyle{fancy}
\fancyhf{}
\setlength{\headheight}{16pt} % Soluciona el warning "headheight is too small"
\lhead{\textbf{IMT-342} | Worksheet Validación URDF}
\rhead{Hito 1 - Semana 6[cite: 7]}

\begin{document}

\section*{A. Identificación del Grupo[cite: 7]}
\textbf{Grupo:} \rule{4cm}{0.15mm} \quad \textbf{Archivo URDF:} \rule{4cm}{0.15mm} \\
\textbf{Integrantes:} \rule{10cm}{0.15mm} \\
\textbf{Base Frame:} \rule{4cm}{0.15mm} \quad \textbf{Tool Frame:} \rule{4cm}{0.15mm}

\section*{B. Validez XML / URDF[cite: 7]}
Herramienta de validación utilizada (Ej. \texttt{check\_urdf}): \rule{5cm}{0.15mm} \\
Resultado: $\square$ Pass \quad $\square$ Fail \qquad Corrección realizada: \rule{5cm}{0.15mm}

\section*{C. Tabla de Definición de Juntas[cite: 7]}
\begin{center}
    \footnotesize
    \renewcommand{\arraystretch}{1.5}
    \noindent\begin{tabularx}{\textwidth}{|c|c|c|c|c|c|X|X|}
    \hline
    \textbf{Joint} & \textbf{Parent} & \textbf{Child} & \textbf{\texttt{xyz}} & \textbf{\texttt{rpy}} & \textbf{\texttt{axis}} & \textbf{Límite Fuente} & \textbf{Justificación Corta} \\ \hline
    J1 & & & & & & & \\ \hline
    J2 & & & & & & & \\ \hline
    J3 & & & & & & & \\ \hline
    J4 & & & & & & & \\ \hline
    J5 & & & & & & & \\ \hline
    J6 & & & & & & & \\ \hline
    \end{tabularx}
\end{center}

\section*{D. Prueba de Movimiento por Junta[cite: 7]}
\begin{center}
    \footnotesize
    \noindent\begin{tabularx}{\textwidth}{|c|X|X|X|X|c|}
    \hline
    \textbf{Joint} & \textbf{Demás fijas en} & \textbf{Valor prueba} & \textbf{Mov. Esperado} & \textbf{Mov. Observado} & \textbf{¿Match?} \\ \hline
    J1 & & & & & $\square$ \\ \hline
    J2 & & & & & $\square$ \\ \hline
    J3 & & & & & $\square$ \\ \hline
    J4 & & & & & $\square$ \\ \hline
    J5 & & & & & $\square$ \\ \hline
    J6 & & & & & $\square$ \\ \hline
    \end{tabularx}
\end{center}

\section*{E. Discrepancia y Corrección (Debugging Real)[cite: 7]}
\textbf{Síntoma detectado:} \rule{12cm}{0.15mm} \\
\textbf{Componente revisado:} $\square$ XML $\square$ Parent/Child $\square$ Origin $\square$ Axis $\square$ Frame conv. \\
\textbf{Corrección y evidencia:} \rule{12cm}{0.15mm} 

\section*{F. Conclusión Final de Validación[cite: 7]}
\textit{Consideramos que el URDF es cinemáticamente defendible porque...} \rule{6cm}{0.15mm} \\
\rule{\textwidth}{0.15mm}

\end{document}
